% Cluster C: certificates and realization (manuscript §10.2)
% Body-only LaTeX; preamble macros and theorem environments per CONVENTIONS.md.

\section{Protected ribbon-routing certificates and realization}
\label{sec:certificates}

The two obstructions of the preceding section calibrate what a realization
theorem can promise.  Theorem~\ref{thm:no-least} shows that even the
two-point package with identity dynamics and no obstacles admits no least
compact connected full carrier, so no canonical minimal choice is available
and a connector must be \emph{supplied}, not selected by an extremal
principle.  Proposition~\ref{prop:access-collision} shows that a monotone,
attachment-complete enlargement can permanently erase a protected prime-end
access even along a conformal, critical-value-free branch: once a set
occupying a protected side has been absorbed, no further monotone
enlargement restores an access through it.  Together these force the shape
of the present section.  The connector data must be recorded \emph{in
advance}, as an explicit finite certificate that names every protected
access, reserves a collar around each one, and routes the connecting tree
through the reserved-free region.  The theorem then owed to the reader is
purely topological: every such certificate is realized by a compact
connected full carrier, the protected accesses survive as prime ends of the
enlarged complement, distinct accesses stay distinct, their cyclic order is
unchanged, and the sides created by the connector itself appear in the
boundary order dictated by the ribbon structure of the tree.  This is the
content of Theorem~\ref{thm:realization} below.  The section closes with
the naturality of the construction under univalent changes of the Jordan
localization (Proposition~\ref{prop:naturality}).  The question of when a
certificate \emph{exists} is deferred to
Section~\ref{sec:criterion}, where it receives an exact answer.

\subsection*{Setup}
Throughout this section $U\subset\C$ is a Jordan disk, that is, a bounded
simply connected domain whose boundary is a Jordan curve.  A
\emph{continuum} is a nonempty compact connected subset of $\C$; a
continuum $K$ is \emph{full} if $\C\setminus K$ is connected, equivalently
if $\Chat\setminus K$ is a simply connected domain containing $\infty$.
Following the source manuscript \cite{MWman}, the full continua to be
joined are called \emph{atoms}.  Let
\[
A_1,\dots,A_s\Subset U
\]
be pairwise disjoint atoms, and let
\[
Z\Subset U\setminus\bigcup_{i}A_i
\]
be a compact set, the frozen \emph{forbidden obstacle and critical-value
set}.  Nothing in this section uses any dynamical structure on the $A_i$ or
on $Z$; the data are purely planar.

Since each $A_i$ is full, $\Omega_i=\Chat\setminus A_i$ is a simply
connected domain, and the prime-end formalism recalled in
Section~\ref{sec:prelim} applies to it \cite{Pom92}.  We use it in the
following normalized form.  A \emph{crosscut} of $\Omega_i$ is an open
Jordan arc $\gamma\subset\Omega_i$ whose closure is a closed arc with two
distinct endpoints on $\partial\Omega_i=\partial A_i$; the crosscut divides
$\Omega_i$ into exactly two subdomains.  A \emph{prime-end chain} (null
chain) is a sequence $(\gamma_n)_{n\ge1}$ of pairwise disjoint crosscuts
whose spherical diameters tend to zero, together with a coherent choice of
complementary domains $D_n$, the \emph{selected sides}: $D_n$ is one of the
two components of $\Omega_i\setminus\gamma_n$, and $D_{n+1}\subset D_n$ for
every $n$.  Two chains are \emph{equivalent} if each eventually divides the
other, that is, every selected domain of either chain contains all
sufficiently deep selected domains of the other; an equivalence class is a
\emph{prime end} of $\Omega_i$.  The record of a chain therefore contains
strictly more than the crosscuts themselves: it contains, for every $n$,
the declaration of which side of $\gamma_n$ is selected.  It is exactly
this declaration that the certificates below protect.

For every protected access to $A_i$, indexed by a finite set of labels $a$,
choose a prime-end crosscut chain
\[
\Gamma_{i,a}=(\gamma_{i,a,n},\,D_{i,a,n})_{n\ge1}
\]
of $\Chat\setminus A_i$.  The chains are truncated, that is, an initial
segment of each is discarded and the indices shifted, so that the terminal
\emph{collar tails}
\[
D_{i,a,1}\setminus A_i
\]
are pairwise disjoint, lie in $U$, and miss $Z$ and the other atoms.  Such
a truncation is possible because the finitely many chains define distinct
accesses with shrinking selected domains, and because $Z$ and the atoms are
compact and disjoint from the impressions being approached.  The set
$D_{i,a,1}$ is called the (terminal, or first) \emph{protected collar} of
the access $(i,a)$; every deeper selected domain $D_{i,a,n}$ lies inside
it.  Figure~\ref{fig:collar} shows a single atom, one protected chain, its
nested selected domains, and its collar.

\begin{figure}[htbp]
\centering
\begin{tikzpicture}[scale=1.05, line join=round]
% ---- selected domains D_1 > D_2 > D_3 (filled before the atom, which
% ---- is drawn on top and cleans up the shared boundary) ----
% D_1: region between crosscut gamma_1 and the oscillating boundary edge
\fill[blue!55, opacity=0.20]
  (0.2,1.1) .. controls (1.5,0.9) and (1.8,-0.9) .. (0.3,-1.2)
  plot [smooth, tension=0.8] coordinates
    {(0.3,-1.2) (-0.5,-0.8) (0.0,-0.4) (-0.6,0.0) (0.1,0.3) (-0.5,0.6) (0.2,1.1)};
% D_2
\fill[blue!55, opacity=0.38]
  (-0.5,0.6) .. controls (0.8,0.45) and (0.95,-0.55) .. (-0.5,-0.8)
  plot [smooth, tension=0.8] coordinates
    {(-0.5,-0.8) (0.0,-0.4) (-0.6,0.0) (0.1,0.3) (-0.5,0.6)};
% D_3
\fill[blue!55, opacity=0.60]
  (0.1,0.3) .. controls (0.55,0.2) and (0.55,-0.25) .. (0.0,-0.4)
  plot [smooth, tension=0.8] coordinates
    {(0.0,-0.4) (-0.6,0.0) (0.1,0.3)};
% ---- crosscuts ----
\draw[very thick, blue!60!black]
  (0.2,1.1) .. controls (1.5,0.9) and (1.8,-0.9) .. (0.3,-1.2);
\draw[very thick, blue!60!black]
  (-0.5,0.6) .. controls (0.8,0.45) and (0.95,-0.55) .. (-0.5,-0.8);
\draw[very thick, blue!60!black]
  (0.1,0.3) .. controls (0.55,0.2) and (0.55,-0.25) .. (0.0,-0.4);
% ---- the atom: an irregular, non-locally-connected-looking blob ----
\fill[gray!25]
  plot [smooth cycle, tension=0.8] coordinates
    {(0.2,1.1) (-1.2,1.6) (-3.0,1.7) (-4.4,0.8) (-4.8,-0.7) (-3.4,-1.9)
     (-1.5,-2.0) (0.3,-1.2) (-0.5,-0.8) (0.0,-0.4) (-0.6,0.0) (0.1,0.3)
     (-0.5,0.6)};
\draw[thick]
  plot [smooth cycle, tension=0.8] coordinates
    {(0.2,1.1) (-1.2,1.6) (-3.0,1.7) (-4.4,0.8) (-4.8,-0.7) (-3.4,-1.9)
     (-1.5,-2.0) (0.3,-1.2) (-0.5,-0.8) (0.0,-0.4) (-0.6,0.0) (0.1,0.3)
     (-0.5,0.6)};
% ---- labels ----
\node at (-2.6,-0.1) {$A_i$};
\node[blue!60!black] at (1.55,0.85) {$\gamma_1$};
\node[blue!60!black] at (1.05,-0.85) {$\gamma_2$};
\node[blue!60!black] at (0.72,0.30) {$\gamma_3$};
\node at (1.02,-0.28) {$D_1$};
\node at (0.52,-0.62) {$D_2$};
\node at (0.30,-0.06) {$D_3$};
% selected-side arrow
\draw[-{Stealth[length=2.5mm]}, thick]
  (2.6,-1.7) node[right, align=left, font=\small]
  {selected side\\ (protected collar $D_1$)}
  .. controls (1.9,-1.3) and (1.6,-0.9) .. (1.15,-0.45);
\end{tikzpicture}
\caption{A protected prime-end crosscut chain at a single atom $A_i$.  The
crosscuts $\gamma_1,\gamma_2,\gamma_3$ are pairwise disjoint arcs of
$\Chat\setminus A_i$ with endpoints on $\partial A_i$ and shrinking
diameters; the shaded regions are the nested selected sides
$D_1\supset D_2\supset D_3$, drawn with increasing opacity toward the
prime end.  The outermost selected domain $D_1$ is the protected collar of
the access: after truncation the collar tail $D_1\setminus A_i$ is
disjoint from every other collar tail, lies in $U$, and misses the
obstacle set $Z$ and the other atoms.  The boundary of the atom is drawn
deliberately oscillating: no local connectivity of $\partial A_i$ is
assumed anywhere in this section, and the access is carried entirely by
the chain, not by any landing curve.}
\label{fig:collar}
\end{figure}

\begin{definition}[Protected ribbon-routing certificate; manuscript
Definition 10.1]\label{def:certificate}
A \emph{protected ribbon-routing certificate} consists of the preceding
data---the Jordan disk $U$, the atoms $A_1,\dots,A_s$, the forbidden set
$Z$, and the truncated protected chains $\Gamma_{i,a}$---together with
pairwise disjoint \emph{free-port arcs} $\alpha_i\colon[0,1]\to U$
satisfying
\[
\alpha_i\cap A_i=\{x_i\},\qquad
\alpha_i((0,1])\cap
\Bigl(Z\cup\bigcup_{j}A_j\cup\bigcup_{j,a}D_{j,a,1}\Bigr)=\varnothing,
\]
where $x_i=\alpha_i(0)$, and an embedded finite tree $G\subset U$ joining
the outer port endpoints $\alpha_1(1),\dots,\alpha_s(1)$ while missing the
same protected and forbidden sets
$Z\cup\bigcup_jA_j\cup\bigcup_{j,a}D_{j,a,1}$.  The planar cyclic order of
half-edges at every vertex of $G$ is retained as part of the certificate
(its \emph{ribbon structure}).  Every access read by a proof must be a
protected chain $\Gamma_{i,a}$ or a declared side of this ribbon tree.
\end{definition}

Given a certificate $R$, define its \emph{carrier}
\begin{equation}\label{eq:C-HR}
H_R \;=\; G\cup\bigcup_i\alpha_i\cup\bigcup_iA_i .
\end{equation}

Four reading notes.  First, the port arc $\alpha_i$ touches its atom at
the single \emph{free port} $x_i$, a point deliberately not lying in any
protected collar: the second displayed condition includes
$a$-indexed collars of the \emph{same} atom $i$, so the port enters the
atom away from every protected access.  Second, ``joining the outer port
endpoints'' is read as part of the embeddedness of the certificate: the
tree $G$ meets each port arc $\alpha_i$ in exactly the point
$\alpha_i(1)$, so that $G\cup\bigcup_i\alpha_i$ is again an embedded
finite tree, with the port arcs attached as pendant edges.  This is the
configuration produced by the construction of
Section~\ref{sec:criterion}, where the routing tree is built in the
complement of the slit port arcs, and it is used in the proof of the
realization theorem below.  Third, the tree $G$ is a
combinatorial object with a planar enhancement: recording the cyclic order
of half-edges at each vertex fixes, up to isotopy, the two sides of every
edge, and it is these \emph{declared ribbon sides} that the realization
theorem converts into new accesses with a prescribed boundary order.
Fourth, the final sentence of the definition is a discipline on all later
arguments, not a property of the data: any subsequent proof that consults
an access of the carrier must consult either a protected chain or a
declared ribbon side, never an incidental access created by the geometry
of a particular embedding.  Figure~\ref{fig:certificate} shows a
certificate with three atoms.

\begin{figure}[htbp]
\centering
\begin{tikzpicture}[scale=1.0, line join=round]
% ---------- Jordan disk U ----------
\draw[very thick] (0,0) ellipse (5.6 and 3.4);
\node at (4.75,2.75) {$U$};
% ---------- obstacle set Z (hatched) ----------
\fill[pattern=north east lines]
  plot [smooth cycle, tension=0.9] coordinates
    {(0.1,1.05) (0.7,1.30) (0.85,1.85) (0.15,2.10) (-0.55,1.75) (-0.45,1.25)};
\draw[thick]
  plot [smooth cycle, tension=0.9] coordinates
    {(0.1,1.05) (0.7,1.30) (0.85,1.85) (0.15,2.10) (-0.55,1.75) (-0.45,1.25)};
\fill[pattern=north east lines]
  plot [smooth cycle, tension=0.9] coordinates
    {(1.0,-2.15) (1.45,-1.95) (1.35,-1.55) (0.85,-1.65)};
\draw[thick]
  plot [smooth cycle, tension=0.9] coordinates
    {(1.0,-2.15) (1.45,-1.95) (1.35,-1.55) (0.85,-1.65)};
\node at (1.35,1.6) {$Z$};
\node at (1.75,-1.75) {$Z$};
% =========== atom A_1 (upper left) ===========
% collar (faces upper-left): shaded selected side + two crosscuts
\fill[blue!55, opacity=0.30]
  (-3.95,1.75) .. controls (-4.55,1.55) and (-4.5,0.95) .. (-3.95,0.85)
  .. controls (-3.6,1.1) and (-3.6,1.5) .. cycle;
\draw[thick, blue!60!black]
  (-3.95,1.75) .. controls (-4.55,1.55) and (-4.5,0.95) .. (-3.95,0.85);
\draw[thick, blue!60!black]
  (-3.85,1.55) .. controls (-4.25,1.45) and (-4.25,1.05) .. (-3.85,1.0);
% blob
\fill[gray!25]
  plot [smooth cycle, tension=0.9] coordinates
    {(-3.4,2.0) (-2.75,1.6) (-2.85,0.95) (-3.5,0.65) (-4.05,0.9) (-3.95,1.5)};
\draw[thick]
  plot [smooth cycle, tension=0.9] coordinates
    {(-3.4,2.0) (-2.75,1.6) (-2.85,0.95) (-3.5,0.65) (-4.05,0.9) (-3.95,1.5)};
\node at (-3.45,1.35) {$A_1$};
\node[blue!60!black, font=\small] at (-4.75,0.75) {$D_{1,a,1}$};
% port
\fill (-2.88,1.0) circle (1.6pt);
\node[font=\small, above right=-2pt] at (-2.92,1.0) {$x_1$};
\draw[very thick, red!70!black]
  (-2.88,1.0) .. controls (-2.35,0.65) and (-2.05,0.35) .. (-1.6,0.05);
\node[red!70!black, font=\small] at (-2.6,0.35) {$\alpha_1$};
% =========== atom A_2 (lower left) ===========
% collar (faces left)
\fill[blue!55, opacity=0.30]
  (-3.75,-1.15) .. controls (-4.45,-1.35) and (-4.45,-2.0) .. (-3.8,-2.2)
  .. controls (-3.5,-1.9) and (-3.5,-1.45) .. cycle;
\draw[thick, blue!60!black]
  (-3.75,-1.15) .. controls (-4.45,-1.35) and (-4.45,-2.0) .. (-3.8,-2.2);
\draw[thick, blue!60!black]
  (-3.65,-1.35) .. controls (-4.15,-1.5) and (-4.15,-1.9) .. (-3.7,-2.0);
% blob
\fill[gray!25]
  plot [smooth cycle, tension=0.9] coordinates
    {(-3.0,-1.1) (-2.35,-1.45) (-2.45,-2.1) (-3.15,-2.35) (-3.8,-2.0) (-3.7,-1.35)};
\draw[thick]
  plot [smooth cycle, tension=0.9] coordinates
    {(-3.0,-1.1) (-2.35,-1.45) (-2.45,-2.1) (-3.15,-2.35) (-3.8,-2.0) (-3.7,-1.35)};
\node at (-3.05,-1.75) {$A_2$};
% port
\fill (-2.42,-1.5) circle (1.6pt);
\node[font=\small, above right=-2pt] at (-2.46,-1.5) {$x_2$};
\draw[very thick, red!70!black]
  (-2.42,-1.5) .. controls (-1.9,-1.15) and (-1.55,-0.95) .. (-1.2,-0.7);
\node[red!70!black, font=\small] at (-1.95,-1.35) {$\alpha_2$};
% =========== atom A_3 (right) ===========
% collar (faces upper right)
\fill[blue!55, opacity=0.30]
  (3.55,1.95) .. controls (4.1,2.2) and (4.65,1.85) .. (4.55,1.3)
  .. controls (4.1,1.35) and (3.75,1.6) .. cycle;
\draw[thick, blue!60!black]
  (3.55,1.95) .. controls (4.1,2.2) and (4.65,1.85) .. (4.55,1.3);
\draw[thick, blue!60!black]
  (3.65,1.75) .. controls (4.0,1.9) and (4.35,1.7) .. (4.3,1.4);
% blob
\fill[gray!25]
  plot [smooth cycle, tension=0.9] coordinates
    {(3.0,1.7) (3.6,1.85) (4.3,1.35) (4.05,0.75) (3.3,0.55) (2.75,1.0)};
\draw[thick]
  plot [smooth cycle, tension=0.9] coordinates
    {(3.0,1.7) (3.6,1.85) (4.3,1.35) (4.05,0.75) (3.3,0.55) (2.75,1.0)};
\node at (3.5,1.2) {$A_3$};
% port
\fill (2.86,0.85) circle (1.6pt);
\node[font=\small, above left=-2pt] at (2.9,0.85) {$x_3$};
\draw[very thick, red!70!black]
  (2.86,0.85) .. controls (2.3,0.55) and (1.95,0.35) .. (1.6,0.1);
\node[red!70!black, font=\small] at (2.5,0.3) {$\alpha_3$};
% =========== ribbon tree G ===========
\fill[red!70!black] (0,-0.45) circle (1.8pt);
\draw[very thick, red!70!black]
  (-1.6,0.05) .. controls (-1.05,-0.15) and (-0.55,-0.35) .. (0,-0.45);
\draw[very thick, red!70!black]
  (-1.2,-0.7) .. controls (-0.75,-0.6) and (-0.4,-0.5) .. (0,-0.45);
\draw[very thick, red!70!black]
  (1.6,0.1) .. controls (1.0,-0.25) and (0.55,-0.4) .. (0,-0.45);
\node[red!70!black] at (0.1,-0.85) {$G$};
% outer endpoints
\fill (-1.6,0.05) circle (1.4pt);
\fill (-1.2,-0.7) circle (1.4pt);
\fill (1.6,0.1) circle (1.4pt);
\end{tikzpicture}
\caption{A protected ribbon-routing certificate with three atoms inside
the Jordan disk $U$.  Each atom $A_i$ carries a protected collar (shaded,
with two nested crosscuts indicated) and a free-port arc $\alpha_i$
(heavy) landing at the single free point $x_i$, chosen away from every
protected collar.  The embedded finite tree $G$ (heavy) joins the outer
port endpoints; the interiors of the $\alpha_i$ and all of $G$ avoid the
hatched obstacle set $Z$, all atoms, and all first collars
$D_{j,a,1}$.  The planar cyclic order of the half-edges at the vertex of
$G$ is part of the certificate.  The carrier is
$H_R=G\cup\bigcup_i\alpha_i\cup\bigcup_iA_i$.}
\label{fig:certificate}
\end{figure}

\subsection*{Fullness of one-point tree attachments}
The first step of the realization is to verify that the carrier
\eqref{eq:C-HR} is a full continuum.  The union
$T=G\cup\bigcup_i\alpha_i$ is itself a finite plane tree: $G$ is a finite
tree and, by the attachment convention recorded after
Definition~\ref{def:certificate}, each $\alpha_i$ meets $G$ in exactly
the point $\alpha_i(1)$, so adjoining the arcs one at a time adds pendant
edges and creates no cycle.  By the port conditions of
Definition~\ref{def:certificate}, $T$ meets each atom $A_i$ in exactly the
point $x_i$ and otherwise misses every atom.  The relevant general fact is
the following lemma, for which we record two proofs: a direct induction on
Janiszewski's separation theorem, and a cohomological argument by
Alexander duality and Mayer--Vietoris.  The first is elementary and
self-contained; the second is the pattern that recurs, for pairwise
disjoint packets, in the packet lemmas of Section~\ref{sec:packets}.

\begin{lemma}[One-point tree attachments are full; manuscript Lemma 10.2]
\label{lem:tree-full}
If a finite plane tree meets each of finitely many pairwise disjoint full
continua in exactly one point and otherwise misses them, then their union
is a full continuum.
\end{lemma}

\begin{proof}[First proof (Janiszewski induction)]
Let $T$ be the finite plane tree and $A_1,\dots,A_s$ the pairwise disjoint
full continua, with $T\cap A_i=\{p_i\}$.  All sets are regarded as compact
subsets of the sphere $\Chat$, and \emph{nonseparating} means that the
complement in $\Chat$ is connected; for a compact subset of $\C$ this is
equivalent to fullness, since the complement in $\Chat$ is open and
connected exactly when the bounded complementary components in $\C$ are
absent.

A finite plane tree is nonseparating.  Indeed, $T$ is compact, connected,
and simply connected, and one may argue by induction on the number of
edges: a single closed arc does not separate the sphere, and if
$T=T'\cup e$ with $e$ a closed pendant edge meeting $T'$ in one endpoint,
then $T'$ and $e$ are nonseparating compacta with connected (one-point)
intersection, so Janiszewski's theorem, quoted below, makes $T$
nonseparating.

We now attach the full continua one at a time.  Put
\[
B_0=T,\qquad B_k=T\cup A_1\cup\dots\cup A_k\quad(1\le k\le s),
\]
so that $B_s$ is the union in question.  Each $B_k$ is compact, and it is
connected because $T$ is connected and each attached $A_j$ meets $T$.  We
claim every $B_k$ is nonseparating.  The case $k=0$ was just proved.
Assume $B_{k-1}$ is nonseparating and consider $B_k=B_{k-1}\cup A_k$.
The intersection is
\[
B_{k-1}\cap A_k
=(T\cap A_k)\cup\bigcup_{j<k}(A_j\cap A_k)
=\{p_k\},
\]
a singleton, hence connected: the atoms are pairwise disjoint and $T$
meets $A_k$ only at $p_k$.  Both $B_{k-1}$ and $A_k$ are compact and
nonseparating.  Janiszewski's separation theorem
\cite{New51,Why42} states that if two closed subsets of the sphere have
connected intersection and neither separates a given pair of points, then
their union does not separate that pair.  Applying this to every pair of
points of $\Chat\setminus B_k$ shows that $B_k$ separates no pair of
points of its complement, that is, $B_k$ is nonseparating.  By induction,
$B_s$ is a compact connected set with connected complement in $\Chat$,
hence a full continuum.
\end{proof}

\begin{proof}[Second proof (Alexander duality and Mayer--Vietoris)]
Work again in $\Chat=S^2$ and retain the notation $B_k$ of the first
proof.  For a compact set $X\subset S^2$, Alexander duality in its
\v{C}ech form \cite{Hat02} gives an isomorphism
\[
\widetilde H_0\bigl(S^2\setminus X;\Z\bigr)\;\cong\;\check H^1(X;\Z),
\]
so $X$ is nonseparating precisely when its first reduced \v{C}ech
cohomology vanishes.  (\v{C}ech cohomology is essential here: the atoms
are arbitrary continua, in general not locally connected, and no
CW-approximation of them is available; \v{C}ech theory has the required
continuity under inverse limits of neighborhoods and satisfies the
Mayer--Vietoris sequence for compact pairs in a normal ambient space.)

The tree $T$ is contractible, so $\check H^1(T;\Z)=0$; and
$\check H^1(A_i;\Z)=0$ for every $i$ because each $A_i$ is full, by the
same duality read in the other direction.  Now induct on $k$.  Suppose
$\check H^1(B_{k-1};\Z)=0$ and write $B_k=B_{k-1}\cup A_k$ with
$B_{k-1}\cap A_k=\{p_k\}$.  The Mayer--Vietoris sequence in reduced
\v{C}ech cohomology for the compact pair $(B_{k-1},A_k)$ reads
\[
\widetilde{\check H}^0(\{p_k\};\Z)
\longrightarrow
\check H^1(B_k;\Z)
\longrightarrow
\check H^1(B_{k-1};\Z)\oplus\check H^1(A_k;\Z).
\]
The left group vanishes because the one-point intersection is connected,
and the right group vanishes by the inductive hypothesis and fullness of
$A_k$.  Exactness forces $\check H^1(B_k;\Z)=0$.  By induction
$\check H^1(B_s;\Z)=0$, and Alexander duality converts this back into
connectedness of $S^2\setminus B_s$.  Since $B_s$ is compact and
connected, it is a full continuum.
\end{proof}

The two proofs establish literally the same induction step, and the
translation between them is the dictionary used throughout this paper:
``nonseparating'' equals ``vanishing first \v{C}ech cohomology,''
``connected intersection'' equals ``vanishing zeroth reduced cohomology of
the overlap.''  The cohomological form is the one that survives when the
finitely many one-point overlaps are replaced by the pairwise
\emph{disjoint} packets of Section~\ref{sec:packets}, where the nerve
argument of Lemma~\ref{lem:packets} takes over.

\subsection*{Survival of protected chains}
The second step is the heart of the realization: enlarging an atom to a
full continuum that stays out of a protected collar does not damage the
protected access.  The chain survives verbatim---the same crosscuts, the
same selected sides---as a prime-end chain of the complement of the
enlarged continuum, and finitely many protected accesses with disjoint
collars stay distinct and keep their cyclic order.

\begin{lemma}[Protected crosscut survival and order; manuscript
Lemma 10.3]\label{lem:crosscut-survival}
Let $H\supset A_i$ be a full continuum and let $(\gamma_n,D_n)_{n\ge1}$ be
one of the protected chains above.  If
\[
(H\setminus A_i)\cap(D_1\setminus A_i)=\varnothing,
\]
then a tail of the same crosscuts is a prime-end chain in
$\Chat\setminus H$.  Its endpoints remain boundary points and its selected
sides contain the old $D_n$.  A finite family of such chains with pairwise
disjoint terminal selected domains induces distinct prime ends of
$\Chat\setminus H$ with unchanged cyclic order.
\end{lemma}

\begin{proof}
Since $H$ is a full continuum, $\Chat\setminus H$ is a simply connected
domain and has its own prime-end theory.  We verify that the tail
$(\gamma_n)_{n\ge2}$, with the same selected sides, is a null chain for
it.

\emph{The selected domains avoid $H$ entirely.}
For every $n$ the selected domain $D_n$ lies in $\Chat\setminus A_i$, so
$D_n\setminus A_i=D_n$, and it is contained in $D_1$ by nestedness.
Hence
\[
H\cap D_n
=(A_i\cap D_n)\cup\bigl((H\setminus A_i)\cap D_n\bigr)
\subset\varnothing\cup\bigl((H\setminus A_i)\cap(D_1\setminus A_i)\bigr)
=\varnothing,
\]
so $D_n\subset\Chat\setminus H$ for every $n$.

\emph{Each deep crosscut remains a crosscut.}
Fix $n\ge2$.  The open arc $\gamma_n$ lies in $D_{n-1}\subset D_1$:
indeed $\gamma_n$ lies in the closure of $D_n\subset D_{n-1}$ and in
$\Chat\setminus A_i$, and the part of $\partial D_{n-1}$ inside
$\Chat\setminus A_i$ is $\gamma_{n-1}$, which $\gamma_n$ misses.  So by
the preceding paragraph the arc $\gamma_n$ misses $H$; its two distinct
endpoints lie on $\partial A_i$, hence in $A_i\subset H$.  Every
neighborhood of either endpoint meets the crosscut interior, which
consists of points of $\Chat\setminus H$, so both endpoints remain
boundary points of the enlarged complement $\Chat\setminus H$, and they
are still distinct.  Thus $\gamma_n$ is again an open Jordan arc in
$\Chat\setminus H$ whose closure is an arc with two distinct endpoints on
$\partial(\Chat\setminus H)$, that is, a crosscut of $\Chat\setminus H$.

\emph{The selected sides are unchanged.}
For $n\ge2$ let $D_n'$ denote the component of
$(\Chat\setminus H)\setminus\gamma_n$ containing $D_n$; this is
well defined because $D_n$ is connected and disjoint from
$H\cup\gamma_n$.  We claim $D_n'=D_n$, so that the containment
$D_n'\supset D_n$ asserted in the statement holds with equality.  The
set $D_n$ is open, and it is relatively closed in
$(\Chat\setminus H)\setminus\gamma_n$: as a component of the open set
$(\Chat\setminus A_i)\setminus\gamma_n$ it satisfies
\[
\partial D_n
\subset\partial A_i\cup\overline{\gamma_n}
\subset A_i\cup\gamma_n
\subset H\cup\gamma_n,
\]
since the closure of the open arc $\gamma_n$ adds only its two
endpoints, which lie in $A_i$; hence no boundary point of $D_n$ belongs
to $(\Chat\setminus H)\setminus\gamma_n$, and every point of that set
lying in $\overline{D_n}$ already lies in $D_n$.  A nonempty connected
subset of an open set that is both open and relatively closed is a
component of it, so $D_n'=D_n$.  Nestedness of the new sides is now
inherited from the original chain:
$D_{n+1}'=D_{n+1}\subset D_n=D_n'$ for every $n\ge2$, and
$\gamma_{n+1}\subset D_n=D_n'$ lies in the selected side of $\gamma_n$,
as a null chain requires.

\emph{The diameters still tend to zero.}
The crosscuts are unchanged as point sets, so their spherical diameters
are the same numbers as before and still tend to zero.  Consequently
$(\gamma_n,D_n')_{n\ge2}$ is a prime-end chain of $\Chat\setminus H$,
which proves the first two assertions.

\emph{Distinctness.}
Now consider a finite family of protected chains satisfying the
hypothesis, with pairwise disjoint terminal selected domains.  Recall
that two null chains of the same domain define the same prime end if and
only if they eventually divide one another: every selected domain of
either chain contains all sufficiently deep selected domains of the
other.  Suppose two chains of the family, with terminal domains
$D_1^{(1)}$ and $D_1^{(2)}$, induced the same prime end of
$\Chat\setminus H$.  Then some deep selected domain of the first
chain---by the preceding step, the same set $D_m^{(1)}$ as before the
enlargement---would be contained in a selected domain of the second
chain, hence in $D_1^{(2)}$; but that deep domain is nonempty and
contained in $D_1^{(1)}$, contradicting
$D_1^{(1)}\cap D_1^{(2)}=\varnothing$.  Disjoint terminal domains
therefore prevent mutual division, and the induced prime ends are
pairwise distinct.  In particular the survival map, from protected
chains to prime ends of $\Chat\setminus H$, is injective on such a
family.

\emph{Cyclic order.}
The prime ends of the simply connected domain $\Chat\setminus H$ carry a
natural cyclic order, the one transported from the unit circle by a
Riemann map \cite{Pom92}.  For a finite family of prime ends represented
by pairwise disjoint crosscuts whose selected sides are also pairwise
disjoint, this order can be read off inside the domain, as follows.
Choose one deep crosscut from each chain, deep enough that the chosen
crosscuts lie in the pairwise disjoint terminal collars; then the chosen
crosscuts are pairwise disjoint and their selected sides are pairwise
disjoint.  Fix three chains of the family, with chosen crosscuts
$c_1,c_2,c_3$ and selected sides $S_1,S_2,S_3$.  Call an embedded triod
$Y$---three arcs with a common branch point, disjoint except at that
point---\emph{computing} if $Y$ lies in the domain and meets the
configuration $\bigcup_{k=1}^{3}(c_k\cup S_k)$ in exactly its three leg
tips, one tip on each $c_k$.  Transporting by a Riemann map turns the
configuration into three pairwise disjoint crosscuts of the unit disk
cutting off pairwise disjoint caps, and a computing triod into an
embedded triod tipping on the three crosscuts and otherwise avoiding
crosscuts and caps.  In the disk, computing triods exist, and the cyclic
order of the three legs at the branch point of any computing triod
agrees with the cyclic order of the three caps on the circle, that is,
with the cyclic order of the three prime ends: the complement of the
three open caps in the closed disk is a closed Jordan region whose
boundary traverses the three crosscuts in the cyclic order of the caps,
and an embedded triod in a Jordan region whose legs end on three
disjoint boundary arcs realizes at its branch point exactly the cyclic
order of those arcs.  Hence, back in the original domain, computing
triods exist for every triple of chains, and each one realizes the
cyclic orientation of its prime-end triple at its branch point.

Now compare the two domains.  By the preceding steps, the chosen
crosscuts and their selected sides are \emph{literally the same sets}
for $\Chat\setminus A_i$ and for $\Chat\setminus H$.  Fix a triple of
chains and choose a computing triod $Y$ in the smaller domain
$\Chat\setminus H$.  Since $\Chat\setminus H\subset\Chat\setminus A_i$
and the configuration is unchanged, $Y$ is simultaneously a computing
triod in $\Chat\setminus A_i$.  Reading the branch point of this one
triod in both domains shows that the prime-end triple has the same
cyclic orientation in $\Chat\setminus H$ as in $\Chat\setminus A_i$.
Every ordered triple therefore retains its orientation, and the cyclic
order of the whole family is unchanged.
\end{proof}

\subsection*{The realization theorem}
Assembling the two lemmas yields the main result of this section.

\begin{theorem}[Protected regulated connector realization; manuscript
Theorem 10.4]\label{thm:realization}
Every protected ribbon-routing certificate $R$ has an actual carrier
$H_R$, given by \eqref{eq:C-HR}, with the following properties:
\begin{enumerate}
\item $H_R\Subset U$ is compact, connected, and full;
\item it contains every $A_i$, misses $Z$, and has singleton attachment
trace $\{x_i\}$ at each atom;
\item every protected crosscut chain remains a prime-end chain in
$\Chat\setminus H_R$;
\item the survival map is injective and preserves cyclic order;
\item the connector-created sides have the boundary order of the retained
ribbon tree; and
\item no local connectivity of any $A_i$ is required.
\end{enumerate}
\end{theorem}

\begin{proof}
\emph{Clause} (1).  As observed before Lemma~\ref{lem:tree-full}, the set
$T=G\cup\bigcup_i\alpha_i$ is a finite plane tree meeting each atom $A_i$
in exactly the point $x_i$ and otherwise missing every atom.
Lemma~\ref{lem:tree-full} applies to $T$ and $A_1,\dots,A_s$ and shows
that $H_R=T\cup\bigcup_iA_i$ is a full continuum.  Each constituent is
compactly contained in $U$, and there are finitely many, so
$H_R\Subset U$.

\emph{Clause} (2).  Containment of the atoms is part of
\eqref{eq:C-HR}.  The certificate conditions of
Definition~\ref{def:certificate} state that the arcs
$\alpha_i((0,1])$ and the tree $G$ miss $Z$, and $Z$ misses every atom by
hypothesis, so $H_R\cap Z=\varnothing$.  The attachment trace at $A_i$ is
the intersection of $A_i$ with the connector: $T\cap A_i=\{x_i\}$ because
$\alpha_i\cap A_i=\{x_i\}$, the other port arcs and $G$ miss $A_i$, and
the atoms are pairwise disjoint.

\emph{Clause} (3).  Fix a protected chain $\Gamma_{i,a}$.  We verify the
hypothesis of Lemma~\ref{lem:crosscut-survival} for $H=H_R$ and the atom
$A_i$.  The difference $H_R\setminus A_i$ is contained in
\[
G\cup\bigcup_j\alpha_j((0,1])\cup\bigcup_{j\ne i}A_j .
\]
The tree $G$ and every half-open port arc $\alpha_j((0,1])$ miss
\emph{all} first collars $D_{j',a',1}$, in particular $D_{i,a,1}$, by
Definition~\ref{def:certificate}; and each atom $A_j$ with $j\ne i$
misses the collar tail $D_{i,a,1}\setminus A_i$ by the truncation
normalization.  Hence
$(H_R\setminus A_i)\cap(D_{i,a,1}\setminus A_i)=\varnothing$: the
connector misses the terminal collar tail, and therefore misses every
later closed collar tail $D_{i,a,n}\subset D_{i,a,1}$ as well.
Lemma~\ref{lem:crosscut-survival} now shows that a tail of
$\Gamma_{i,a}$ is a prime-end chain in $\Chat\setminus H_R$.

\emph{Clause} (4).  The truncation normalization makes the terminal
selected domains of distinct protected chains pairwise disjoint.  For
each atom $A_i$, the final part of Lemma~\ref{lem:crosscut-survival}
with $H=H_R$ gives distinctness and preservation of cyclic order for the
subfamily of chains protecting $A_i$.  The distinctness argument of the
lemma uses only disjointness of the terminal selected domains and
survival of both chains in $\Chat\setminus H_R$, so it applies verbatim
to a pair of chains protecting different atoms as well.  The survival
map is therefore injective on the full family over all $(i,a)$, and it
preserves the cyclic order of every protected family with a common
source domain.

\emph{Clause} (5).  Let $N$ be a small regular neighborhood of the finite
connector tree $T$ in the plane, chosen so thin that $N$ misses $Z$ and
meets each atom only near its port point.  A regular neighborhood of a finite plane tree is a closed
disk, and its boundary circle traverses every edge of $T$ on two sides:
following $\partial N$, one passes along each edge once in each
direction, and at a vertex of valence $d$ the boundary visits the $d$
half-edges in exactly their planar cyclic order.  The resulting cyclic
\emph{boundary word}, in which every edge appears twice with opposite
orientations, is precisely the boundary word of the ribbon structure
retained in the certificate.  The accesses to $H_R$ created by the
connector---the two sides of every edge of $G$ and of every port
arc---are read along this boundary circle, so their boundary order is the
boundary order of the retained ribbon tree, as claimed.

\emph{Clause} (6).  The argument uses crosscut chains only: atoms enter
through their prime-end chains and through the single point $x_i$, never
through a boundary parametrization.  No Euclidean boundary extension of
any $A_i$---in particular no continuity of a Riemann map up to the
boundary, equivalently no local connectivity of any atom---is assumed at
any step.
\end{proof}

\begin{remark}\label{rem:C-no-lc}
Clause (6) is the point of the whole construction and deserves emphasis.
Realization arguments for connectors frequently assume that accesses
land, that is, that radial limits of Riemann maps exist at the relevant
boundary points; by Carath\'eodory's theorem this amounts to local
connectivity of the atom boundaries.  The certificate formalism removes
that hypothesis: an access is a chain of crosscuts with declared sides,
an object available for an arbitrary full continuum, and
Lemma~\ref{lem:crosscut-survival} shows it is preserved verbatim under
any enlargement respecting the collar.  This is what allows the framework
to be applied, in Section~\ref{sec:decoration}, to atoms that are not
known---and are not claimed---to be locally connected.
\end{remark}

\subsection*{Univalent naturality}
The certificate data and the carrier construction are natural under
univalent change of the Jordan localization.

\begin{proposition}[Univalent naturality; manuscript Proposition 10.5]
\label{prop:naturality}
If $\phi\colon U\to V$ is an orientation-preserving conformal
homeomorphism of Jordan disks and every certificate entry is transported
by $\phi$, then
\begin{equation}\label{eq:C-naturality}
H_{\phi_*R}=\phi(H_R).
\end{equation}
The protected survival map and the ribbon boundary order commute with the
prime-end map induced by $\phi$.
\end{proposition}

\begin{proof}
Since $U$ and $V$ are Jordan disks, Carath\'eodory's extension theorem
\cite{Pom92} extends $\phi$ to a homeomorphism
$\overline U\to\overline V$ of the closed disks.  By the Schoenflies
theorem the restricted boundary homeomorphism
$\partial U\to\partial V$ extends to a homeomorphism of the complementary
closed disks $\Chat\setminus U\to\Chat\setminus V$, and gluing the two
extensions along the boundary curves produces an orientation-preserving
homeomorphism $\Phi\colon\Chat\to\Chat$ with $\Phi|_U=\phi$.

A sphere homeomorphism preserves every notion entering
Definition~\ref{def:certificate} and Theorem~\ref{thm:realization}: it
carries compact connected sets to compact connected sets and complements
to complements, hence preserves fullness; it carries open Jordan arcs
with endpoints on a boundary to the same, hence crosscuts to crosscuts;
it maps the two complementary sides of a crosscut to the two sides of the
image crosscut, hence selected sides to selected sides; it carries null
chains to null chains (spherical diameters are not literally preserved,
but by uniform continuity of $\Phi$ the image crosscuts of a null chain
again have diameters tending to zero); it preserves intersections, so the
port conditions $\alpha_i\cap A_i=\{x_i\}$ and all disjointness
requirements transport; and, being orientation preserving, it preserves
the planar cyclic order of half-edges at the tree vertices and the cyclic
order of prime ends.  Consequently the transported data---the disk $V$,
the atoms $\phi(A_i)$, the obstacle $\phi(Z)$, the chains
$\Phi(\Gamma_{i,a})$, the ports $\phi\circ\alpha_i$, and the tree
$\phi(G)$ with its transported ribbon structure---form a protected
ribbon-routing certificate in $V$.  The defining union \eqref{eq:C-HR}
then transports term by term:
\[
H_{\phi_*R}
=\phi(G)\cup\bigcup_i\phi(\alpha_i)\cup\bigcup_i\phi(A_i)
=\phi\Bigl(G\cup\bigcup_i\alpha_i\cup\bigcup_iA_i\Bigr)
=\phi(H_R),
\]
which is \eqref{eq:C-naturality}.  Because $\Phi$ maps
$\Chat\setminus H_R$ onto $\Chat\setminus H_{\phi_*R}$ carrying null
chains to null chains and preserving their division relation, it induces
a bijection of prime ends under which the survival map of
Lemma~\ref{lem:crosscut-survival} for $H_R$ corresponds to the survival
map for $H_{\phi_*R}$, and the ribbon boundary words of
Theorem~\ref{thm:realization}(5) correspond letter by letter.
\end{proof}

Consequently one actual anchor certificate transports through the
complete degree-one corridor of the source manuscript \cite{MWman}.  We
record the manuscript's caveat: this transport does \emph{not} prove
equality of the transported carrier with any carrier copy supplied by the
$\psi$-ql immersion.  Naturality moves a certificate and its realization
along a univalent map; it does not identify the result with an
independently constructed carrier, and no such identification is claimed
anywhere in this paper.

With realization settled, the burden shifts to existence: when do the
port arcs and the routing tree demanded by
Definition~\ref{def:certificate} exist at all?
Section~\ref{sec:criterion} answers this question exactly, in terms of a
single component of the frozen free domain $\Omfree$ admitting
simultaneous singleton ports to every required atom.
